\subsection{局部可读而全局受阻：Čech 障碍、规范差与曲率候选}\label{subsec:physical-spacetime-skeleton-local-global-obstruction}

上一小节说明，只要局部观测族足够分离，时间迹可以在算子层面重建空间传播核的谱。然而，这个结论仍然只是一种局部—半全局结论：它告诉我们，许多空间信息可以通过局部时间演化恢复出来；但它并没有保证这些局部可读信息一定能无缝拼成一个单一的全局对象。换言之，局部可读并不自动推出全局可拼接。本小节的目的，就是把这种局部上已经成立、但整体上仍受阻的现象形式化出来，并把它解释为一种 Čech 型障碍、规范差与曲率候选。

这一步是必要的，因为前文已经把空间解释成共同支撑—局部比较—区域共识的最小涌现结构。既然整体是由局部相容族拼出来的，那么整体失败就不会首先表现为一个数值错误，而会表现为：某些局部截面在两两重叠区上虽然可比较，却无法在三重重叠乃至更高重叠上满足一致性闭合。因此，本小节并不新增新的对象层原语，而是把局部 forcing、局部读出与局部拼接推到它们自然会出现的边界：局部上看似都成立，整体上却不存在一个单一全球代表。

以下固定上一小节得到的涌现区域系统 \(\mathcal X:=\mathcal X_{\mathrm{em}}(\mathcal O)\)，并在其覆盖结构上讨论局部读出、规范过渡与全局化成败。

\subsubsection{读出预层与局部可读性}

\paragraph{定义 15.9.1（读出预层）}
固定一个对象 \(a\)（地址、事件标签、模式标签或其他可指称对象）。对每个区域 \(U\in\mathcal X\)，定义
\[
\mathscr R_a(U)
\]
为对象 \(a\) 在区域 \(U\) 上的一切合法局部读出、局部截面或局部代表所成之集合。若 \(V\subseteq U\)，给定限制映射
\[
\rho^U_V:\mathscr R_a(U)\to \mathscr R_a(V).
\]
并要求：
\begin{enumerate}
  \item \(\rho^U_U=\mathrm{id}\)；
  \item 若 \(W\subseteq V\subseteq U\)，则
  \[
  \rho^V_W\circ \rho^U_V=\rho^U_W.
  \]
\end{enumerate}
于是 \(U\mapsto \mathscr R_a(U)\) 构成对象 \(a\) 的一个读出预层。

这里的读出必须理解得尽可能宽：它可以是第 \ref{sec:recursive-addressing} 节里的局部截面，可以是第 \ref{sec:emergent-arithmetic} 节里的局部正规形，也可以是第 \ref{sec:pom} 节投影语法中的局部词、局部门或局部证书。关键不在其具体类型，而在于：它们都 obey 限制与局部比较。

\paragraph{定义 15.9.2（局部可读与全局可读）}
设 \(U\in\mathcal X\)，覆盖
\[
\mathfrak U=\{U_i\to U\}_{i\in I}\in \operatorname{Cov}(U).
\]
\begin{enumerate}
  \item 若对每个 \(i\in I\) 都有
  \[
  \mathscr R_a(U_i)\neq\varnothing,
  \]
  则称对象 \(a\) 在覆盖 \(\mathfrak U\) 下局部可读；
  \item 若
  \[
  \mathscr R_a(U)\neq\varnothing,
  \]
  则称对象 \(a\) 在区域 \(U\) 上全局可读。
\end{enumerate}
于是，局部可读但全局受阻就意味着：某个对象在每个局部块上都有合法读出，但整个区域上没有一个统一读出。

\subsubsection{严格拼接与规范拼接}

局部读出要想拼成整体，至少有两种不同强度的要求。最强的是严格一致，较弱的是允许规范差。

\paragraph{定义 15.9.3（严格兼容族）}
在覆盖 \(\mathfrak U=\{U_i\to U\}_{i\in I}\) 下，若选定
\[
s_i\in \mathscr R_a(U_i),
\]
并满足对任意 \(i,j\) 都有
\[
\rho^{U_i}_{U_i\cap U_j}(s_i)=\rho^{U_j}_{U_i\cap U_j}(s_j),
\]
则称 \((s_i)_{i\in I}\) 是一个严格兼容族。

若存在 \(s\in \mathscr R_a(U)\) 使
\[
\rho^U_{U_i}(s)=s_i\qquad(\forall i),
\]
则称 \(s\) 是该兼容族的一个严格全局拼接。

这就是通常意义上的层条件。若 \(\mathscr R_a\) 本身是层，那么严格兼容就足以推出唯一全局拼接。

但在许多技术场景中，局部代表并不会真的在重叠区上一字不差地相等，而只是相差一个允许的规范变换。因此还要定义规范拼接。

\paragraph{定义 15.9.4（规范群预层与规范兼容族）}
设 \(\mathscr G_a\) 是作用在 \(\mathscr R_a\) 上的一个群预层，即对每个区域 \(U\)，\(\mathscr G_a(U)\) 是一群，且若 \(V\subseteq U\)，存在限制同态
\[
\mathscr G_a(U)\to \mathscr G_a(V).
\]

称局部读出族 \((s_i)_{i\in I}\) 是一个规范兼容族，若对每个重叠区 \(U_{ij}:=U_i\cap U_j\)，都存在
\[
g_{ij}\in \mathscr G_a(U_{ij})
\]
使
\[
s_i|_{U_{ij}}=g_{ij}\cdot s_j|_{U_{ij}}.
\]

这里的 \(g_{ij}\) 就是局部代表之间的规范差。严格兼容对应于所有 \(g_{ij}=e\)。

\subsubsection{Čech 障碍与三重重叠闭合}

一旦允许规范差，真正关键的就不再是两两重叠，而是三重重叠上的闭合。

\paragraph{定义 15.9.5（三重重叠障碍元）}
对三重重叠
\[
U_{ijk}:=U_i\cap U_j\cap U_k,
\]
定义
\[
c_{ijk}:=g_{ij}g_{jk}g_{ki}\in \mathscr G_a(U_{ijk}).
\]
称 \((c_{ijk})\) 为规范兼容族的三重重叠障碍元。

若对所有 \(i,j,k\) 都有
\[
c_{ijk}=e,
\]
则称该规范兼容族满足三重闭合条件。

这正是 Čech 型障碍最原始的形式：局部代表两两之间都能用规范差联系起来，但这些规范差绕一圈回来未必闭合为单位元。

\paragraph{命题 15.9.6（全局规范拼接的必要条件）}
若局部读出族 \((s_i)\) 存在一个全局代表 \(s\in \mathscr R_a(U)\)，以及局部规范 \(h_i\in \mathscr G_a(U_i)\)，使
\[
s_i=h_i\cdot s|_{U_i},
\]
则其三重重叠障碍元满足
\[
c_{ijk}=e\qquad(\forall i,j,k).
\]

\paragraph{证明}
在重叠区上有
\[
s_i|_{U_{ij}} = h_i|_{U_{ij}}\cdot s|_{U_{ij}},
\qquad
s_j|_{U_{ij}} = h_j|_{U_{ij}}\cdot s|_{U_{ij}}.
\]
故
\[
g_{ij}=h_i h_j^{-1}
\quad\text{于 }U_{ij}\text{ 上}.
\]
于是
\[
c_{ijk}
=
g_{ij}g_{jk}g_{ki}
=(h_i h_j^{-1})(h_j h_k^{-1})(h_k h_i^{-1})
=e.
\]
证毕。

这说明：若三重障碍元不为单位，就不可能存在一个统一的全局代表把所有局部代表重新校正为同一对象。

\paragraph{推论 15.9.7（Čech 障碍阻断全局可读）}
若对象 \(a\) 在覆盖 \(\mathfrak U\) 下局部可读，且存在一组局部代表 \((s_i)\) 的三重障碍元满足
\[
c_{ijk}\neq e
\]
于某个三重重叠区上成立，则 \(a\) 在区域 \(U\) 上不能全局可读，至少不能以该规范群 \(\mathscr G_a\) 的意义全局拼接。

\subsubsection{规范差、holonomy 与全局扭曲}

即使三重重叠障碍元全部消失，也不代表一定能得到一个全局选定代表。因为此时仍可能存在非平凡的 \(1\)-余循环类：局部上能粘，但全局只能粘成一个扭曲对象，而不是单一的平凡化代表。

\paragraph{定义 15.9.8（规范差 \(1\)-余循环）}
若规范兼容族满足三重闭合条件 \(c_{ijk}=e\)，则 \((g_{ij})\) 构成一个 Čech \(1\)-余循环。若存在局部规范 \(h_i\in\mathscr G_a(U_i)\) 使
\[
g_{ij}=h_i h_j^{-1},
\]
则称该 \(1\)-余循环是平凡的；否则称其为非平凡规范差。

非平凡规范差表示：局部代表之间虽然两两与三重重叠都兼容，但仍不存在一个全局单值代表可把它们同时拉平。

这时，全局受阻不再表现为局部不一致，而表现为全局扭曲。其最直接的物理读法，就是 holonomy：沿一个闭合重叠链回到起点后，局部代表虽然仍然合法，但获得了一个不可去掉的整体相位或群元素。

若选定一条闭合重叠链
\[
U_{i_0}\to U_{i_1}\to \cdots \to U_{i_m}=U_{i_0},
\]
则定义其 holonomy 为
\[
\mathrm{Hol}(\gamma):=
g_{i_0 i_1}g_{i_1 i_2}\cdots g_{i_{m-1}i_m}.
\]
若 \(\mathrm{Hol}(\gamma)\neq e\)，则表明局部表示在该闭环上具有不可去除的全局扭曲。

因此，本小节的全局受阻至少有两层：第一，三重重叠障碍不消失，连局部—局部—局部的一致性都无法成立；第二，三重重叠已闭合，但 \(1\)-余循环类仍非平凡，局部可拼，整体却只有扭曲对象而无单一平凡代表。

\begin{proposition}[局域骨架先于全局平凡化闭包]\label{prop:physical-spacetime-local-skeleton-before-global-trivialization}
沿用本小节记号，并设前几小节已在某覆盖
\[
\mathfrak U=\{U_i\to U\}_{i\in I}
\]
上分别构造出观察者、时间、空间与因果的局域骨架读出 \(s_i\in\mathscr R_a(U_i)\)。
则有：
\begin{enumerate}
  \item \textbf{局域构造层独立存在：}这些 \(s_i\) 已经给出一个合法的局域骨架族，其成立不依赖后续是否能全局平凡化。
  \item \textbf{全局平凡化的必要条件：}若存在一个单值的全局骨架 \(s\in\mathscr R_a(U)\) 及局部规范 \(h_i\in\mathscr G_a(U_i)\) 使
  \[
  s_i=h_i\cdot s|_{U_i},
  \]
  则对所有三重重叠必有 \(c_{ijk}=e\)，并且 \(1\)-余循环 \((g_{ij})\) 必为 Čech 平凡。
  \item \textbf{障碍层阻断闭包升级：}因此，只要存在非平凡三重障碍元或非平凡 holonomy，局域骨架就不能被提升为后续全局 untwisted skeleton；在此之前，不得调用定理 \ref{thm:physical-spacetime-admissible-einstein-closure-main} 所要求的 admissible 全局 Einstein 闭包。
\end{enumerate}
\end{proposition}

% --- Distillation writeback ---
\begin{proposition}[有限一维骨架上的 reduced holonomy 平凡化门控]\label{prop:distill-gromov-physical-finite-reduced-holonomy-promotion-gate}
\textbf{依赖状态：rigidity.}
令 \(\Gamma=(V,E)\) 为有限连通无向图，并为每条无向边选取两个相反的有向边记号 \(e,\bar e\)。设 \(G\) 为群，且给定边标记
\[
g_e\in G,\qquad g_{\bar e}=g_e^{-1}.
\]
对有向路径 \(\gamma=e_1\cdots e_r\)，采用约定
\[
H_g(\gamma):=g_{e_1}g_{e_2}\cdots g_{e_r}.
\]
取一棵生成树 \(T\subset\Gamma\) 与根顶点 \(o\)。对每条不属于 \(T\) 的有向边 \(e=(u,v)\)，令
\[
\ell_e:=P_{o,u}\,e\,P_{v,o}
\]
其中 \(P_{a,b}\) 表示树 \(T\) 中从 \(a\) 到 \(b\) 的唯一有向路径。则以下三条等价：
\begin{enumerate}
  \item 对所有 \(e\notin T\)，有 \(H_g(\ell_e)=1_G\)；
  \item 对任意闭合有向路径 \(\gamma\)，有 \(H_g(\gamma)=1_G\)；
  \item 存在顶点规范 \(a_v\in G\ (v\in V)\)，使每条有向边 \(e=(u,v)\) 满足
  \[
  g_e=a_u^{-1}a_v.
  \]
\end{enumerate}
因此，在本小节的局部读出覆盖上，若边过渡先经某个商群或 reduced register 后得到这样的 \(g_e\)，则只需核验生成树外的基本闭路 holonomy；全部通过时，reduced 边过渡可被提升为全局平凡化数据。若某个基本闭路未通过，则它给出阻断该 reduced 平凡化的有限 holonomy 见证。
\end{proposition}

% --- Distillation writeback ---
\begin{proposition}[约化 holonomy 通过后的核残差审计]\label{distill:gromov-formal_versus_genuine_realization_families-physical-reduced-kernel-holonomy-audit}
\textbf{依赖状态：conditional audit.}
令 \(\Gamma=(V,E)\) 为有限连通无向图，并为每条无向边选取两个相反的有向边记号 \(e,\bar e\)。令
\[
\kappa:G\twoheadrightarrow \overline G
\]
为交换群之间的满同态，并记 \(K:=\ker\kappa\)。设给定 \(G\)-值边过渡
\[
g_e\in G,
\qquad
g_{\bar e}=-g_e.
\]
记 \(\bar g_e:=\kappa(g_e)\)。假设约化 holonomy 审计已经通过，即存在顶点规范 \(\bar a_v\in\overline G\) 使每条有向边 \(e=(u,v)\) 满足
\[
\bar g_e=\bar a_v-\bar a_u.
\]
任取顶点 lift
\[
a_v\in G,
\qquad
\kappa(a_v)=\bar a_v,
\]
并定义核残差边账本
\[
k_e:=g_e-(a_v-a_u)\in G
\qquad(e=(u,v)).
\]
则：
\begin{enumerate}
  \item 每个 \(k_e\) 实际属于 \(K\)，且 \(k_{\bar e}=-k_e\)。对任意闭合有向路径 \(\gamma=e_1\cdots e_m\)，有
  \[
  \sum_{r=1}^m g_{e_r}=\sum_{r=1}^m k_{e_r}\in K.
  \]
  \item 原边过渡存在 \(G\)-值全局平凡化，即存在 \(b_v\in G\) 使每条有向边 \(e=(u,v)\) 满足
  \[
  g_e=b_v-b_u,
  \]
  当且仅当核残差 \((k_e)\) 是 \(K\)-值余边界，即存在 \(c_v\in K\) 使
  \[
  k_e=c_v-c_u\qquad(e=(u,v)).
  \]
  等价地，取任意生成树 \(T\subset\Gamma\) 后，只需核验所有树外边给出的基本闭路上的核残差和均为零。
  \item 因此，约化 holonomy 审计只证明了边过渡在商群 \(\overline G\) 中可平凡化；若某个基本闭路具有非零核残差和，则它是阻断完整 \(G\)-值全局平凡化的有限 holonomy 见证，必须进入本小节意义下的离散障碍账本，而不能被约化通过结论吸收。
\end{enumerate}
\end{proposition}
\begin{proof}
先证第一项。若 \(e=(u,v)\)，则
\[
\kappa(k_e)=\kappa(g_e)-\kappa(a_v-a_u)=\bar g_e-(\bar a_v-
\bar a_u)=0,
\]
故 \(k_e\in K\)。反向边满足
\[
k_{\bar e}=g_{\bar e}-(a_u-a_v)=-g_e+a_v-a_u=-k_e.
\]
若闭合路径写作 \(e_r=(v_{r-1},v_r)\)，且 \(v_m=v_0\)，则
\[
\sum_{r=1}^m g_{e_r}
=\sum_{r=1}^m \bigl((a_{v_r}-a_{v_{r-1}})+k_{e_r}\bigr)
=\sum_{r=1}^m k_{e_r},
\]
因为顶点项沿闭路望远镜相消。右端属于 \(K\)，第一项得证。

再证第二项。若存在 \(b_v\in G\) 使 \(g_e=b_v-b_u\)，令
\[
d_v:=b_v-a_v.
\]
对边 \(e=(u,v)\)，有
\[
k_e=g_e-(a_v-a_u)=(b_v-b_u)-(a_v-a_u)=d_v-d_u.
\]
由于
\[
\kappa(d_v)-\kappa(d_u)=\kappa(k_e)=0,
\]
并且 \(\Gamma\) 连通，\(\kappa(d_v)\) 与 \(v\) 无关。固定根顶点 \(o\)，令
\[
c_v:=d_v-d_o.
\]
则 \(c_v\in K\)，且
\[
k_e=c_v-c_u.
\]
故核残差为 \(K\)-值余边界。

反过来，若存在 \(c_v\in K\) 使 \(k_e=c_v-c_u\)，定义
\[
b_v:=a_v+c_v.
\]
则对每条有向边 \(e=(u,v)\)，
\[
b_v-b_u=(a_v-a_u)+(c_v-c_u)=(a_v-a_u)+k_e=g_e.
\]
于是原边过渡在 \(G\) 中全局平凡化。

最后说明基本闭路核验。对任意 \(K\)-值边账本 \((k_e)\)，若它是余边界，则任意闭合路径上的和显然为零。反之，取生成树 \(T\) 与根 \(o\)，定义
\[
c_v:=\sum_{e\in P_{o,v}} k_e,
\]
其中 \(P_{o,v}\) 为树中从 \(o\) 到 \(v\) 的唯一有向路径。对树边，立刻有 \(k_{(u,v)}=c_v-c_u\)。若 \(e=(u,v)\notin T\)，其基本闭路为 \(P_{o,u}eP_{v,o}\)。该闭路核残差和为零等价于
\[
c_u+k_e-c_v=0,
\]
即 \(k_e=c_v-c_u\)。故所有树外基本闭路通过当且仅当 \((k_e)\) 为余边界。

第三项只是上述等价性的审计解释：商群中的平凡化不能自动消去 \(K\)-方向的闭路残差；完整全局平凡化必须再核验核残差类为零。证毕。
\end{proof}

% --- End distillation writeback ---

\begin{proof}
先证 (3) 推 (2)。若 \(\gamma=e_1\cdots e_r\) 为闭合路径，写 \(e_j=(v_{j-1},v_j)\) 且 \(v_r=v_0\)。由 (3) 得
\[
H_g(\gamma)=a_{v_0}^{-1}a_{v_1}a_{v_1}^{-1}a_{v_2}\cdots a_{v_{r-1}}^{-1}a_{v_r}=1_G.
\]
于是 (2) 推 (1) 显然成立。

再证 (1) 推 (3)。对每个顶点 \(v\)，定义
\[
a_v:=H_g(P_{o,v}).
\]
若 \(e=(u,v)\) 是树边，则树路径满足 \(P_{o,v}=P_{o,u}e\) 或等价地在相反方向使用 \(g_{\bar e}=g_e^{-1}\)，故总有
\[
a_v=a_u g_e,\qquad\text{即}\qquad g_e=a_u^{-1}a_v.
\]
若 \(e=(u,v)\) 不在树中，则由假设
\[
1_G=H_g(P_{o,u}eP_{v,o})=a_u g_e a_v^{-1},
\]
从而同样得到 \(g_e=a_u^{-1}a_v\)。故 (3) 成立。证毕。
\end{proof}

% --- End distillation writeback ---

\begin{proof}
\textbf{(1)} 这只是重述“局部可读”的定义：每个 \(U_i\) 上的 \(s_i\) 都是合法局部代表，因此局域骨架已被构造出来。

\textbf{(2)} 三重障碍元消失由命题 15.9.6 直接给出。
若存在全局骨架 \(s\) 与局部规范 \(h_i\)，则在每个重叠区上
\[
g_{ij}=h_i h_j^{-1},
\]
故定义 15.9.8 意义下的 Čech \(1\)-余循环是平凡的。

\textbf{(3)} 由 (2) 的逆否命题，只要 \(c_{ijk}\neq e\) 或 \((g_{ij})\) 非平凡，就不存在单值的全局 untwisted skeleton。
而定理 \ref{thm:physical-spacetime-admissible-einstein-closure-main} 的 admissible 闭包要求的正是这种已完成全局化的骨架对象；
故在障碍未消去前，最多只能停留在局域骨架或扭曲对象层，不能越级宣称全局 Einstein 闭包。
证毕。
\end{proof}
\leanverified{paper\_physical\_spacetime\_local\_skeleton\_before\_global\_trivialization}

\subsubsection{曲率候选：从离散障碍到连续近似}

当重叠系统足够细，且规范群具有局部线性化或可交换化近似时，上述 Čech 障碍可进一步被解释为一种离散曲率。

\paragraph{定义 15.9.9（离散曲率候选）}
设 \(\mathscr G_a\) 在某个局部极限中可嵌入某个李群，且对每个重叠区上的规范差可写成
\[
g_{ij}=\exp(A_{ij}),
\]
其中 \(A_{ij}\) 取值于某个线性空间 \(\mathfrak g\)。定义三角胞腔 \((i,j,k)\) 上的离散曲率候选为
\[
\Omega_{ijk}:=
A_{ij}+A_{jk}+A_{ki}.
\]

若群近似可交换，则
\[
c_{ijk}=e
\quad\Longleftrightarrow\quad
\Omega_{ijk}=0.
\]

因此，\(\Omega_{ijk}\) 正是 Čech 三重障碍的线性化版本。

这说明：曲率在这里并不是从连续联络预设出来的，而是从局部代表能否在三重重叠上闭合自然导出的。换言之，本章中的曲率首先是一个拼接失败量，然后才在细极限中逼近通常的几何曲率。

若覆盖继续细化，且 \(\Omega_{ijk}\) 对细化满足一致的归并规律，则可把它看作某个连续 \(2\)-形式 \(F\) 的离散近似。于是：\(F=0\) 对应局部数据可无扭曲地全局拼接；\(F\neq 0\) 对应存在真实的全局几何或规范障碍。这就是曲率候选这个说法的意义：在本章中，曲率不是起点，而是由 Čech 障碍在细化极限中导出的二次对象。

\subsubsection{与前文各层的关系}

这一节并不是突然把外来的层论或规范论硬塞进来，而是前文 forcing 链条在全局失败处的自然延伸。

与第 \ref{sec:recursive-addressing} 节的关系是最直接的：第 \ref{sec:recursive-addressing} 节中的局部截面与 \(\mathsf{NULL}\) 已经表明，一个对象可以在某个局部区域上可读，而在更大的区域上仍不进入决定包络。本小节给出了这种现象的全局形式化版本：局部可读但全局受阻，不一定只是没有全局截面，还可能是只存在扭曲全局对象而无平凡化代表。

与第 \ref{sec:emergent-arithmetic} 节的关系也很自然：第 \ref{sec:emergent-arithmetic} 节中的规范横截面与值保持重写，默认是在一个已经可被全局选定的值纤维上进行。本小节则说明：若全局读出本身带有非平凡 Čech 障碍或 holonomy，则所谓唯一规范横截面只能局部定义，不能无条件全局定义。也就是说，值层正规形的全球存在性，隐含依赖于本小节的拼接无障碍条件。

与小节 \ref{subsec:physical-spacetime-skeleton-space-shared-support} 至小节 \ref{subsec:physical-spacetime-skeleton-time-space-exchange} 的关系则在于：共同支撑、局部共识、资源几何与时间迹/空间谱重建，都默认有一个至少在某个层面上可以全局组织的局部系统。而本小节指出，这种组织本身可能失败，失败的最小单位正是 Čech 障碍。于是，本小节为后续任何全局谱、全局热核或全局压力之类对象划出真正边界：它们只有在局部读出系统可全局拼接，或至少扭曲可控时，才是真正意义上的全局不变量。

\subsubsection{本节结论}

本小节的核心结论可以压缩为四句话。第一，局部可读不推出全局可读。第二，两两重叠上的兼容不推出三重重叠上的闭合。第三，三重闭合即便成立，也未必推出全局平凡代表存在。第四，全局受阻的最小度量不是数值误差，而是 Čech 障碍、规范差与其离散曲率候选。

因此，本小节真正完成的，是把空间的最小涌现从小节 \ref{subsec:physical-spacetime-skeleton-space-shared-support} 推进到了全局空间何以会失败的层面：不是所有局部共识都能无缝提升为全局几何，而一旦失败，这种失败不是噪声，而是一个可被记录、可被分类、可被线性化的结构量。

也正是在这个意义上，本小节为后面继续讨论谱与资源提供了真正的边界条件：谱不是装饰，而是建立在可拼接闭环之上的边界不变量；资源不是附会，而是消除障碍、维持共识、压低 holonomy 与曲率所必须付出的代价。
